\documentclass{beamer}
\usetheme{boxes}
\usecolortheme{default}
\usepackage[T1]{fontenc}
\usepackage{amsmath,amssymb}
\usepackage{csquotes}

% Bibliographic metadata was checked against C:/mytex/svozil.bib.
% References are deliberately rendered only as descriptive DOI hyperlinks.
\newcommand{\doiref}[2]{%
  \href{https://doi.org/#1}{\textcolor{blue!65!black}{#2}}}
\newcommand{\sourcefoot}[1]{%
  \par\vfill{\tiny\color{black!60}DOI sources: #1\par}}

% Title information from C:/mytex/2026-Krakow-pres.tex.
\title{Indeterminism and Randomness\\ in Quantum and Classical Physics}
\author{\texorpdfstring{Karl Svozil\\ITP, TU Wien, Vienna, Austria}%
{Karl Svozil, ITP, TU Wien, Vienna, Austria}}
\date{Workshop on Algorithmic Randomness, Krakow, Poland\\
2026-08-31--2026-09-01\\
\url{https://svozil.github.io/publications/2026-Krakow-pres.pdf}}

\begin{document}

\begin{frame}
\titlepage
\end{frame}

\begin{frame}{The Principle of Sufficient Reason (PSR)}

\begin{block}{Leibniz's formulation (\emph{Monadology} {\S}32, 1714)}
\small
\enquote{\textit{And that of sufficient reason, in virtue of which we hold
that no fact can be real or existent, no statement true, unless there be a
sufficient reason why it is so and not otherwise, although these reasons
usually cannot be known by us.}}
\vskip2pt
\hfill\footnotesize --- G.\,W.\ Leibniz, \emph{Monadologie}, {\S}32
(written in French; publ.\ posthumously)
\end{block}

\vspace{-0.2em}
\begin{exampleblock}{How Leibniz meant it}
\small
\begin{itemize}
  \item \textbf{Metaphysical, not formal:} an axiom about \emph{reality} and
        its intelligibility --- never stated in a calculus or derived from
        logic.
  \item \textbf{\enquote{Reason} = ground/explanation}, wider than
        \emph{cause}: it answers \emph{why so and not otherwise}.
  \item \textbf{Epistemically modest:} the sufficient reason often involves an
        infinite analysis --- fully surveyable only by God.
\end{itemize}
\end{exampleblock}

\end{frame}

\begin{frame}{Suppose the answer is ``not real''---\\
or inaccessible by the available means}
\small
\begin{enumerate}
  \setlength{\itemsep}{0.55em}
  \item \textbf{A Viennese passer-by:}
        \enquote{Where in Vienna is St.~Paul's Cathedral?}
        The named cathedral has no Viennese referent; the famous one is in
        London. Vienna's central cathedral is St.~Stephen's.
  \item \textbf{A car:}
        \enquote{How much milk is in your refrigerator?}
        An ordinary car has no relevant sensor or data channel.
  \item \textbf{A large language model:}
        Ask for a fact that is neither represented in its learned evidence nor
        supplied in the context or through tools.
  \item \textbf{A quantum system:}
        Prepare normalized $\lvert\psi\rangle$ and ask the yes/no question
        $P_\phi=\lvert\phi\rangle\!\langle\phi\rvert$, with
        $0<\lvert\langle\phi\vert\psi\rangle\rvert^2<1$.
        The prepared state is not an eigenstate of this query.
\end{enumerate}
\end{frame}

\begin{frame}{What happens when an answer is enforced?}
\small
\begin{enumerate}
  \setlength{\itemsep}{0.55em}
  \item \textbf{The passer-by} may repair the false premise, ask for
        clarification, substitute St.~Stephen's, or simply guess.
  \item \textbf{The car} can return \enquote{unavailable}, request access to a
        refrigerator sensor, or emit a value unsupported by its data.
  \item \textbf{The LLM} may halucinate a plausible but ungrounded answer.
        Alignment and tools can encourage abstention or evidence-seeking, but
        cannot guarantee either.
  \item \textbf{The quantum measurement} records $1$ with
        $p=\lvert\langle\phi\vert\psi\rangle\rvert^2$ and $0$ with $1-p$.
        In the complementary qubit case, $p=1/2$.
\end{enumerate}

\begin{alertblock}{The contrast}
\footnotesize
The first three cases concern false premises, missing access, or unsupported
inference. The fourth is a positive Born-rule prediction: an admissible outcome
must occur, although the prepared state fixes only its probability.
\end{alertblock}
\end{frame}

\begin{frame}{Stace: realism exceeds possible experience}
\footnotesize
\begin{block}{The claim he targets}
Realism says that some entities sometimes exist without being experienced by
any finite mind. Stace does \emph{not} prove their nonexistence. His conclusion
is epistemic: their unexperienced existence is neither known nor made probable.
\end{block}

\begin{columns}[T,onlytextwidth]
\begin{column}{0.485\textwidth}
\begin{exampleblock}{Three routes fail}
\textbf{Perception:} experiencing the unexperienced is contradictory.
\par\smallskip
\textbf{Induction:} no unexperienced existence has ever been observed.
\par\smallskip
\textbf{Deduction:} an object's absence during an observational gap is not
logically inconsistent.
\end{exampleblock}
\end{column}

\begin{column}{0.485\textwidth}
\begin{exampleblock}{Two replies fail}
\textbf{Egocentric predicament:} ignorance does not prove nonexistence, but it
provides no reason for existence either.
\par\smallskip
\textbf{Causal continuity:} it presupposes that unobserved processes and laws
persist --- the very point to be proved.
\end{exampleblock}
\end{column}
\end{columns}

\begin{alertblock}{Stace's conclusion}
Persistence through \enquote{interperceptual intervals} is a simplifying
mental construction or fiction, not an established observer-independent fact.
\end{alertblock}

% svozil.bib key: stace.
\sourcefoot{\doiref{10.1093/mind/XLIII.170.145}{Stace 1934}.}
\end{frame}

\begin{frame}{Quantum ``anti-realism'': localized indefiniteness}
\footnotesize
\begin{block}{Theorem 2: a finite witness for every incompatible query}
Let $n\geq3$ and prepare a pure state $\lvert\psi\rangle$, so the eigenstate
assumption fixes $v(P_\psi)=1$. For every unit $\lvert\phi\rangle$ with
$0<\lvert\langle\psi\vert\phi\rangle\rvert<1$, one can effectively construct a
finite set $O_\phi$ of rank-one projectors containing $P_\psi$ and $P_\phi$ for
which no admissible noncontextual value assignment can make $P_\phi$ definite.
\end{block}

\begin{exampleblock}{Angle from the prepared ray}
Let $\theta=\arccos\lvert\langle\psi\vert\phi\rangle\rvert$, with
$0\leq\theta\leq\pi/2$.
\begin{description}
  \setlength{\itemsep}{0.3em}
  \item[$\theta=0$:] same ray, $P_\phi=P_\psi$; definite value $1$.
  \item[$0<\theta<\pi/2$:] neither aligned nor orthogonal; $P_\phi$ is value
        indefinite in the theorem's sense.
  \item[$\theta=\pi/2$:] orthogonal ray; definite value $0$ under admissibility.
\end{description}
\end{exampleblock}

\begin{alertblock}{What ``anti-realism'' means here}
Almost every rank-one proposition lacks a pre-existing \emph{noncontextual}
$0/1$ value relative to the prepared state. The result assumes the eigenstate
condition and admissibility; each target $P_\phi$ has its own finite witness
$O_\phi$. It does not rule out contextual hidden-variable accounts or every
form of metaphysical realism.
\end{alertblock}

% svozil.bib key: 2015-AnalyticKS.
\sourcefoot{\doiref{10.1063/1.4931658}{Abbott--Calude--Svozil 2015, Theorem 2}.}
\end{frame}

\begin{frame}{Quantum indeterminacy in Zeilinger's account}
\small
\begin{block}{Finite information}
Zeilinger's foundational principle says that an \emph{elementary} system
represents the truth value of one proposition --- one bit of information.
\end{block}

\begin{exampleblock}{A complementary query}
For a yes/no proposition not fixed by that bit, quantum mechanics supplies
Born probabilities rather than a pre-existing eigenvalue. Zeilinger interprets
the individual outcome as irreducibly random: the \enquote{message of the
quantum}.
\end{exampleblock}

\begin{alertblock}{PSR, Gaps and continuous creation}
Forget about The Principle of Sufficient Reason---its false!
Philipp Frank argued that in such cases mechanical laws may leave \emph{gaps} (for \emph{miracles}): an observed
initial state need not determine a unique observable future. By analogy,
putatively uncaused event actualization can be called \emph{``creatio continua''}.
In theology the term usually means ongoing creation or sustenance, not a
physical randomness mechanism.
\end{alertblock}

% svozil.bib keys: zeil-99; zeil-05_nature_ofQuantum; franke;
% svozil-2016-quantum-hokus-pokus.
\sourcefoot{\doiref{10.1023/A:1018820410908}{Zeilinger 1999};
\doiref{10.1038/438743a}{Zeilinger 2005};
\doiref{10.1007/978-94-011-5516-8}{Frank 1998};
\doiref{10.3354/esep00171}{Svozil 2016}.}
\end{frame}

\begin{frame}[shrink=3]{Quantum examples proposed for \emph{creatio continua}}
\footnotesize
\begin{columns}[T,onlytextwidth]
\begin{column}{0.485\textwidth}
\begin{block}{Candidate events}
\begin{itemize}
  \setlength{\itemsep}{0.5em}
  \item \textbf{Spontaneous emission:} the registered decay time and emitted
        mode are probabilistic, although the closed atom--field state can
        evolve unitarily.
  \item \textbf{Stimulated emission:} an external field drives the transition;
        it can be coherent and unitary, so it is not by itself a collapse.
  \item \textbf{Query mismatch:} measuring $P_\phi$ on a non-eigenstate
        $\lvert\psi\rangle$ produces a Born-distributed result.
\end{itemize}
\end{block}
\end{column}

\begin{column}{0.485\textwidth}
\begin{block}{Measurement dynamics}
\begin{itemize}
  \setlength{\itemsep}{0.5em}
  \item \textbf{von Neumann Process 1:} projection/collapse --- a nonunitary
        state update.
  \item \textbf{Process 2:} unitary, reversible evolution preserving inner
        products and distances.
  \item \textbf{Partial trace:} a reduced subsystem description, not a
        separate objective collapse.
\end{itemize}
\end{block}
\end{column}
\end{columns}

\begin{alertblock}{Ways to cope?}
\begin{itemize}
  \setlength{\itemsep}{0.5em}
  \item
Everett retains universal Process~2 dynamics.
Thus \enquote{creatio continua} belongs to collapse-based readings, not to Everett's fundamental
dynamics.
\item
Infinite succession of Wigner's Friends, modelled by infinite tensor products, yields unitary non-equivalence.
\item Lemmas by L\'evy's and Johnson--Lindenstrauss on FAPP-orthogonality.
\end{itemize}
\end{alertblock}

% svozil.bib keys: svozil-pac; v-neumann-49; everett.
\sourcefoot{\doiref{10.1007/978-3-319-70815-7_14}{Svozil, ``Vacuum Fluctuations''};
\doiref{10.1007/978-3-642-61409-5}{von Neumann};
\doiref{10.1103/RevModPhys.29.454}{Everett 1957};
\href{http://www.numdam.org/item/CM_1939__6__1_0/}{\textcolor{blue!65!black}{von Neumann 1939}};
\doiref{10.1007/s10701-025-00903-9}{Svozil 2026}.}
\end{frame}


\begin{frame}{Kreisel and Specker: limits of computation}
\footnotesize
\begin{block}{Kreisel's criterion}
A theory is \emph{mechanistic} if every observable sequence or real number
it defines is recursive in the data that determine it. This computability
criterion cuts across deterministic and probabilistic theories.
\end{block}

\begin{exampleblock}{A recursive curve with no computable maximizer}
There is a recursively continuous $f\colon[0,1]\to\mathbb{R}$, with a
recursive modulus of continuity, whose maximum is attained at no recursive
$x$. The maximum value is the computable number $1/36$; its maximizing
arguments are noncomputable.
\end{exampleblock}

\begin{alertblock}{A Specker sequence: computable terms, noncomputable limit}
There is a computable, monotonically increasing, bounded sequence $(q_n)$ of
rational numbers whose limit $\alpha=\lim_n q_n$ is not a computable real.
Consequently, the convergence has no computable rate: no computable $N(m)$
can guarantee $|q_n-\alpha|<2^{-m}$ for every $n\geq N(m)$. Every finite
term can be generated effectively; the asymptotic value cannot.
\end{alertblock}

% svozil.bib keys: kreisel; specker49; specker57.
\sourcefoot{\doiref{10.1007/BF00484949}{Kreisel 1974};
\doiref{10.2307/2267043}{Specker 1949};
\doiref{10.1007/978-3-0348-9259-9_12}{Specker 1959 (1990 reprint)}.}
\end{frame}

\begin{frame}{Classical indeterminism and randomness}
\footnotesize
\begin{columns}[T,onlytextwidth]
\begin{column}{0.485\textwidth}
\begin{block}{Continuum and chaos}
\begin{itemize}
  \setlength{\itemsep}{0.32em}
  \item $\mathbb{R}$ is uncountable.
  \item Lebesgue-almost every real is Martin-L\"of random; choice alone
        supplies neither sampling nor randomness.
  \item Chaos can amplify small uncertainty; injectivity is unnecessary.
\end{itemize}
\end{block}
\end{column}

\begin{column}{0.485\textwidth}
\begin{exampleblock}{Beyond ordinary computation}
\begin{itemize}
  \setlength{\itemsep}{0.45em}
  \item The diagonal Ackermann function is total and computable, but not
        primitive recursive.
  \item Chaitin's $\Omega_U$ is a machine-dependent, uncomputable,
        Martin-L\"of-random \emph{real} --- not an evolution function.
  \item The busy-beaver function $\Sigma(n)$ is total and uncomputable; it
        eventually exceeds every total computable function.
  \item A partial $f$ is undefined on some inputs: that is failure or
        nontermination, not stochastic choice.
\end{itemize}
\end{exampleblock}
\end{column}
\end{columns}

% svozil.bib keys: MartinLof1966; ackermann:28 (DOI added here);
% s00032-006-0064-2; rado.
\sourcefoot{\doiref{10.1016/S0019-9958(66)80018-9}{Martin-L\"of 1966};
\doiref{10.1007/BF01459088}{Ackermann 1928};
\doiref{10.1007/s00032-006-0064-2}{Chaitin 2007};
\doiref{10.1002/j.1538-7305.1962.tb00480.x}{Rad\'o 1962}.}
\end{frame}

\begin{frame}{Vacuous ``laws'' and correlations through Ramsey's theorem}
\footnotesize
\begin{block}{Ramsey's theorem}
Fix positive integers $k,r,c$.  For all sufficiently large $N$, every
$c$-colouring of the $k$-element subsets of an $N$-element set contains an
$r$-element subset whose $k$-element subsets all have the same colour.
\end{block}

\begin{exampleblock}{Some regularity is unavoidable}
Encode a finitely classified relation among $k$ observations as a colour.
In a sufficiently large data set, one can then select a homogeneous
substructure: within that selected set, the same relation holds throughout.
Thus an apparently striking correlation may be guaranteed to occur
\emph{somewhere}, even when no dynamics or causal mechanism produced it.
\end{exampleblock}

\begin{alertblock}{When does a ``law'' become vacuous?}
If the variables, coding, scale, and subset are chosen only after inspecting
the data, Ramsey regularity can be redescribed as a ``law.''  Physical content
requires pre-specified selection criteria, predictions beyond the fitted
cases, and robustness under new observations.  Ramsey's theorem guarantees
existence of a pattern---not its causal significance.
\end{alertblock}

\sourcefoot{\doiref{10.1112/plms/s2-30.1.264}{Ramsey 1930}.}
\end{frame}


\begin{frame}{Martin-L\"of randomness and prefix-free incompressibility}
\footnotesize
\begin{columns}[T,onlytextwidth]
\begin{column}{0.485\textwidth}
\begin{block}{All effective statistical tests}
A Martin-L\"of-random infinite bit sequence passes every algorithmically
implementable test of randomness whose rejection criterion captures only an
effectively vanishing fraction of fair-coin sequences.  No such test singles
it out as exceptional.
\end{block}
\end{column}

\begin{column}{0.485\textwidth}
\begin{exampleblock}{Algorithmic incompressibility}
Descriptions are \emph{self-delimiting}, or prefix-free, programs: no complete
program begins another.  For each initial segment of $n$ bits, the
\emph{shortest} program producing it is essentially $n$ bits long; its saving
is bounded by one constant independent of $n$.  Longer descriptions always
exist.
\end{exampleblock}
\end{column}
\end{columns}

\begin{alertblock}{Levin--Schnorr theorem}
The two viewpoints select exactly the same infinite sequences: Martin-L\"of
randomness is equivalent to prefix-free incompressibility.  \emph{Algorithmic
randomness} is a broader umbrella term; in this talk it means Martin-L\"of
randomness.  No finite prefix can certify this infinite-sequence property.
\end{alertblock}

% svozil.bib keys: MartinLof1966; ch:75; Nies2009.
\sourcefoot{\doiref{10.1016/S0019-9958(66)80018-9}{Martin-L\"of 1966};
\doiref{10.1145/321892.321894}{Chaitin 1975};
\doiref{10.1093/acprof:oso/9780199230761.001.0001}{Nies 2009}.}
\end{frame}


\begin{frame}{Random sequences: almost surely, not necessarily}
\footnotesize
\begin{block}{Single outcomes do not determine a sequence law}
Value indefiniteness or irreducible chance concerns what fixes each outcome
$X_n$.  By itself, it does not specify how successive outcomes are related.
\end{block}

\begin{exampleblock}{The additional process postulate}
Assume independent repetitions governed by an effectively calculable
probability law.  For $n$ fair bits,
$\Pr(X_1\cdots X_ns)=2^{-n}$.  Under this model, the infinite outcome sequence
is Martin-L\"of random with probability one.
\end{exampleblock}

\begin{alertblock}{Probability one $\neq$ logical necessity}
This is an \emph{almost-sure} theorem, not a guarantee for every conceivable
infinite realization.  Algorithmic randomness follows almost surely from the
additional process model---not from single-event indeterminacy alone.  No
finite prefix can certify Martin-L\"of randomness.
\end{alertblock}

% svozil.bib keys: 2012-incomput-proofsCJ; MartinLof1966.
\sourcefoot{\doiref{10.1103/PhysRevA.86.062109}{Abbott et al. 2012};
\doiref{10.1016/S0019-9958(66)80018-9}{Martin-L\"of 1966}.}
\end{frame}

\begin{frame}{Further reading}
\centering
\vfill
{\large Karl Svozil\par}
\vspace{0.65em}
{\LARGE\itshape
\doiref{10.1007/978-3-319-70815-7}{Physical (A)Causality}\par}
\vspace{0.45em}
{\large Determinism, Randomness and Uncaused Events\par}
\vspace{0.9em}
{\normalsize Fundamental Theories of Physics, Vol.~192\\
Springer, 2018 --- open access\par}
\vspace{1em}
\begin{minipage}{0.78\textwidth}
\centering
A fuller treatment of provable unknowns, quantum and classical
indeterminism, and the status of putatively uncaused events.
\end{minipage}
\vfill
{\Large \color{blue}Thank you for your attention\par}
\vspace{0.5em}
% svozil.bib key: svozil-pac.
\end{frame}

\end{document}
